\gls{ELF} is the standard executable format under Linux operating systems.
The work done in this thesis heavily involves the \gls{ELF} structure, with particular focus on some sections 
used for relocation, alongside the process of executing an \gls{ELF}.
\subsection{ELF Structure}
    In this subsection the process of \gls{ELF} execution is discussed, with a focus on the parts of the execution flow 
    that will influence this work.\\
    The \textbf{ELF Header} is a structure located at the beginning of the \gls{ELF} and contains some important information 
    about the type of \gls{ELF}, the target operating system \dots.
    Alongside this information there is the location of two important data structures inside the \gls{ELF}:
    \begin{itemize}
        \item The \textbf{Program Header Table}, which defines the \textit{Execution View}
        \item The \textbf{Section Header Table}, which defines the \textit{Linking View}
    \end{itemize}
    The \textbf{Program Header Table} contains information needed to create a runtime image of the program in memory, via 
    the \textbf{Segments} \textit{(i.e. which parts of the file must be loaded into memory, which shared 
    objects are needed \dots)}.
    Each segment contains the access flags for that part of memory (i.e. Read/Write/Execute).
    The \textbf{Section Header Table} contains information about the \textbf{Sections} that compose the binary, 
    together with their mapping to the segments (e.g. which segment the \textit{.text} section is mapped to).
    Among all the sections, some contain specific structures used to perform specific tasks:
    \begin{itemize}
        \item \textit{\textbf{.dynamic}: The dynamic section contains information used by the Dynamic Linker 
            (e.g. the dependencies, the interpreter location \dots)}.
        \item \textit{\textbf{.dynsym}: List of symbols defined and used by the binary. 
            If the binary is not stripped, there is also the \textbf{.symtab} section, which contains a more complete list of symbols}.
        \item \textit{\textbf{.hash or .gnu.hash}: Structures used for efficient symbol lookup during symbol resolution}\footnote{
            The \textbf{.gnu.hash} section is the more efficient version of \textbf{.hash}.
        }.
        \item \textit{\textbf{.rela.dyn and .rela.plt}: Both are relocation tables. The first is the standard one, 
        the second is used in the case of \textbf{Lazy Binding}}.
        \item \textit{\textbf{.shstrtab, .dynstr and .strtab}: Hold lists of strings referenced by other parts of the \gls{ELF}.
            The first holds section names, the second holds the names of the symbols contained in the \textbf{.dynsym} section, 
            and the last one holds the names of the entire \textbf{.symtab} section}.
        \item \textit{\textbf{.got and .got.plt}: Contain the \gls{GOT}, filled by the dynamic linker with the 
            addresses of external functions and variables. The first is the standard one, the second is used by the \gls{PLT}}.
        \item \textit{\textbf{.plt}: Contains the \gls{PLT}, which is an executable section made up of trampolines. In particular, it 
            contains stubs for functions that might be located in another \gls{ELF} or \gls{SO}. 
            The stub references the corresponding \gls{GOT} entry containing the resolved address of the function.
            This section is used in the case of \textbf{Lazy Binding}}.
        \item \textit{\textbf{.init and .init\_array}: Sections used for initialization. \textbf{.init} contains initialization code, 
            while \textbf{.init\_array} contains the addresses of initialization functions}.
        \item \textit{\textbf{.fini and .fini\_array}: Sections used for termination handling. 
            \textbf{.fini} contains termination code, 
            while \textbf{.fini\_array} contains the addresses of termination functions}.
    \end{itemize}
\subsection{ELF Linking}
    Before actually digging deep into \gls{ELF} loading, the linking process must be reviewed.
    The \textbf{Linking phase} consists of joining together different pieces of code belonging to different files.
    For example, if a program uses the \texttt{printf} function, it needs to take the code of printf from the libc library.
    This process of joining together the executable that uses printf with the printf code located in another file 
    is called \textbf{Linking}. Linking can be of two different types: 
    \begin{itemize}
        \item \textbf{Static linking}: The needed code (e.g. the instructions of printf) is copied inside the binary 
        at \textit{Link Time}. This causes the size of the binary to be bigger (because the code 
        of external functions actually lies inside it) and consequently wastes RAM space (if 10 binaries all use the 
        printf function, the code of printf is copied inside all 10 binaries).
        \item \textbf{Dynamic linking}: The needed code is not copied inside the binary; instead, the binary contains information (i.e. 
        the address of the external function or variable) used to locate the needed resource.
        The resolution happens at runtime, when the Dynamic Linker finds and loads the library into the process's virtual address space and 
        resolves the references to the symbols. This allows different processes to share the same \gls{SO}.
    \end{itemize}
    Since nowadays the most commonly used approach is \textbf{Dynamic Linking}, the following discussion will focus on it.\\
    As stated before, the linking process can be used for both external variables and external functions.
    In both cases the process is the same, since both functions and variables can be exported as symbols.\\
    When dealing with functions there is the possibility of using \textbf{Lazy Binding}.
    This technique consists of resolving the symbol only during the first invocation of the function, rather than at load time 
    as happens with variables.
    In order to achieve \textbf{Lazy Binding} a stub section is used (i.e. \textbf{.plt}).
    When a program uses a library function, a \gls{PLT} entry is added to the .plt section.
    \begin{figure}[ht]
    \centering
            \begin{subfigure}[h]{0.7\textwidth}
            \centering
            \includegraphics[width=\linewidth]{LazyBindingExample.png}
            \caption{An example of puts function resolution using Dynamic Linking and Lazy Binding before the actual resolution}
            \label{fig:lazy_resolution_plt_1}
        \end{subfigure}
        \vspace{0.5cm}
        \begin{subfigure}[h]{0.7\textwidth}
            \centering
            \includegraphics[width=\linewidth]{LazyBindingExampleAfterResolution.png}
            \caption{Same binary as \cref{fig:lazy_resolution_plt_1} but with symbol resolved}
            \label{fig:lazy_resolution_plt_2}
        \end{subfigure}
        \caption{An example of Lazy Binding involving .plt and .got}
        \label{fig:lazy_resolution}
    \end{figure}
    As can be seen in \cref{fig:lazy_resolution}, the relocation is resolved lazily upon the first invocation of the function. 
    When control is transferred to the corresponding \gls{PLT} entry, the entry performs an indirect jump through 
    its associated \gls{GOT} entry. Before the symbol is resolved, this \gls{GOT} entry contains an address that 
    redirects execution back to the \gls{PLT}, eventually reaching the \gls{PLT} resolver stub. 
    The resolver then invokes the dynamic linker, which performs the symbol lookup, resolves the relocation, 
    and updates the \gls{GOT} entry with the resolved address of the function, as can be seen in \cref{fig:lazy_resolution_plt_2}. 
    Subsequent calls can therefore jump directly to the resolved address through the \gls{GOT}.

\subsection{ELF Loading}
    The next useful step in order to understand the chosen approach consists of going into detail on the ELF Loading phase.
    There are 3 main actors that allow a binary to be executed properly:
    \begin{itemize}
        \item \textbf{Kernel}: responsible for the actual \texttt{execve} invocation and 
            allocation of the different memory segments. Moreover, it is the one that executes 
            the dynamic linker used to resolve dependencies prior to the execution of the binary.
        \item \textbf{Dynamic Linker}: Loads all the needed libraries into the process memory, 
            resolves the external dependencies, prepares the .got and .plt, and executes the initializers of the libraries.
        \item \textbf{Target Binary}: Invoked by the Dynamic Linker, it prepares the environment for the main function and, 
            in the end, runs it.
    \end{itemize}

    \begin{figure}[ht]
        \centering
        \includegraphics[width=300pt]{dynamicallyLinkedProgramLoad.png}
        \caption{Dynamically linked ELF loading procedure}
        \label{fig:dynamicallyLinkedProgramLoad}
    \end{figure}
    All these procedures can be seen in detail in \cref{fig:dynamicallyLinkedProgramLoad}.
    The dynamic linker part is the one on which I will focus in this dissertation.
    The dynamic linker:
    \begin{itemize}
        \item \textit{Loads all the DT\_NEEDED entries}: these are the dependencies (i.e. the \gls{SO}s needed by the binary).
            For example, if the binary uses an external function contained in \textit{libc.so}, the library from which the 
            function is exported must be added as a dependency.
        \item \textit{Resolves the relocations}: a focus on relocations will be presented in the next subsection.
        \item \textit{Arranges the \gls{GOT} and \gls{PLT}}.
        \item \textit{Executes all the initializer functions (constructors) for each dependency}: For each \textit{DT\_NEEDED} 
            entry found by the dynamic linker, it executes all the functions contained in the \textit{.init\_array} section of 
            that entry. These are functions declared with \texttt{\_\_attribute\_\_((constructor))}.
        \item \textit{Finds and executes the binary entry point}.
    \end{itemize}

\subsection{ELF Symbols and Relocations}
    The dynamic linker uses symbols to link different ELFs together.
    \begin{itemize}
        \item A component (an \gls{ELF} or a \gls{SO}) can export a symbol representing a variable or a function inside it.
        \item Another component can reference the exported symbol and use it (if the first component is a dependency of the second).
        \item The linker acts as the intermediary between the two components, responsible for matching the symbol names 
            (the exported one and the referenced one) and storing the address of the exported symbol in the second component's address space 
            in such a way that this component can access it.
    \end{itemize}
    All the symbols belonging to an \gls{ELF} (both imported and exported) are stored in the \textit{.dynsym} section.
    \begin{figure}[ht]
        \centering
        \includegraphics[width=\linewidth]{dynsym_example.png}
        \caption{Example of a \textit{.dynsym} table for a shared object}
        \label{fig:dynsym_example}
    \end{figure}
    \begin{figure}[ht]
        \centering
        \includegraphics[width=\linewidth]{reladynexample.png}
        \caption{Example of a \textit{.rela.dyn} relocation table extracted from a binary}
        \label{fig:reladyn_example}
    \end{figure}
    In \cref{fig:dynsym_example} a list of symbols extracted from a \gls{SO} can be seen.
    The first 4 symbols are external symbols that must be resolved. This can be seen because the Value column is 0 and Ndx is UND (undefined).
    The last 7, instead, are symbols declared inside the shared object.
    When the visibility is set to \textit{DEFAULT}, this means that the symbol can be visible outside the shared object.

    The process of resolving symbol references is called \textbf{Relocation}. The relocation algorithm is shown in 
    \cref{alg:relocation_resolution}.
    Relocating a symbol means resolving the symbol (i.e. finding the address of that symbol in the original object), and writing 
    that address inside the binary that has the relocation.
    Specifically, the symbol address is written inside the \gls{GOT}.
    In \cref{fig:reladyn_example}, the address at which the symbol address must be written can be seen in the ``Offset'' column.
    \begin{algorithm}
        \caption{Relocation resolution at ELF load time}
        \label{alg:relocation_resolution}
        \begin{algorithmic}[1]
            \REQUIRE $L = [obj_1,\dots,obj_n]$ \COMMENT{The objects that have a dependency (DT\_NEEDED) must already be loaded in memory}
            \ENSURE for each $obj \in L$, the respective GOT entry contains the resolved runtime address of the external symbol.
            
            \FOR{$obj \in L$}
                \STATE $Reloc\_table \gets obj.\texttt{GET\_SECTION(".rela.dyn")}$
                \FOR{$reloc \in Reloc\_table$}
                    \IF{$reloc.\text{type} = \texttt{R\_*\_RELATIVE}$}
                        \STATE $addr \gets obj.\texttt{loadBias} + r.\text{addend}$ 
                        \COMMENT{\textcolor{green!60!black}{The symbol does not need to be resolved. The address must be adjusted in order to enable ASLR}}
                    \ELSE
                        \IF{$reloc.\text{type} \in \{\texttt{R\_*\_GLOB\_DAT} , \texttt{R\_*\_JUMP\_SLOT}\}$} 
                            \STATE $sIdx \gets reloc.\text{symbol\_index}$
                            \STATE $symbolName \gets obj.\texttt{dynstr}[\, obj.\texttt{dynsym}[sIdx].\text{name} \,]$ 
                            \COMMENT{\textcolor{green!60!black}{The name of the symbol is read from \textit{.dynsym}}}
                            \STATE $target \gets \textsc{ResolveSymbol}(symbolName, L, obj)$
                            \IF{$target = \text{NULL}$}
                                \STATE \textbf{error}: undefined symbol $symbolName$ 
                            \ELSE
                                \STATE $(defObj, defIdx) \gets target$ \COMMENT{\small \textcolor{green!60!black}{\textit{defObj} is the object that exports the symbol, \textit{defIdx} is the index in the .dynsym table of \textit{defObj}}}
                                \STATE $addr \gets defObj.\texttt{loadBias} + defObj.\texttt{dynsym}[defIdx].\text{value}$
                                \COMMENT{\small \textcolor{green!60!black}{The address of the resolved symbol is computed as base\_address (of the object that exports the symbol) plus the offset at which the symbol is stored}}
                            \ENDIF
                        \ENDIF
                    \ENDIF
                    \STATE write $addr$ in $obj.\texttt{got}$ at offset $reloc.\text{offset}$ \COMMENT{\small \textcolor{green!60!black}{The destination address at which to write the resolved symbol address is specified in the relocation entry}}
                \ENDFOR
            \ENDFOR
        \end{algorithmic}
        
    \end{algorithm}
    \begin{algorithm}
        \caption{\textsc{ResolveSymbol} function used by \cref{alg:relocation_resolution}}
        \label{alg:symbol_resolution}
        \begin{algorithmic}[1]
            \REQUIRE name $name$, object list $L$, object requesting the resolution $requester$
            \ENSURE pair $(obj, idx)$ of the first symbol found, or NULL if the symbol has not been defined by any object.
            \FOR{$obj \in L$ \COMMENT{\small \textcolor{green!60!black}{The order is: first the executable, then all the DT\_NEEDED dependencies in order of definition}}}
                \STATE $idx \gets \textsc{GnuHashLookup}(name,\ obj.\texttt{gnuHash},\ obj.\texttt{dynsym},\ obj.\texttt{dynstr})$
                \COMMENT{\small \textcolor{green!60!black}{As seen before, the hashes are used in order to speed up the symbol lookup}}
                \IF{$idx \neq \text{NOT\_FOUND}$ \textbf{\&\&} $obj.\texttt{dynsym}[idx].\text{st\_shndx} \neq \texttt{SHN\_UNDEF}$}
                    \STATE \textbf{return} $(obj, idx)$ 
                \ENDIF
            \ENDFOR
            \STATE \textbf{return} NULL
        \end{algorithmic}
    \end{algorithm}